\begin{center}
	\large{\textbf{KATA PENGANTAR}}
\end{center}
\indent Puji Tuhan. Segala hormat dan syukur dipanjatkan ke hadirat Tuhan Yang Maha Esa karena telah menaruh berkat, petunjuk, kekuatan, dan tuntunan jalan kepada penulis, sehingga penulis dapat menyelesaikan Tugas Akhir yang berjudul:
\begin{center}
	\textbf{``GRAVITASI EINSTEIN DAN EINSTEIN-MAXWELL PADA 2+1 DIMENSI DALAM TEORI GAUGE CHERN-SIMONS''}
\end{center}
sebagai salah satu syarat kelulusan Program Sarjana Fisika FSAD ITS.\\
\indent Dalam penyusunan dan penulisan Tugas Akhir ini, penulis mendapatkan banyak dukungan serta bantuan dari berbagai pihak. Oleh karena itu, penulis menyampaikan rasa terima kasih kepada:
\begin{enumerate}[label=\arabic*.]
	\item Teruntuk diri sendiri, terima kasih atas segala kegigihan, ketabahan, dan keberanian untuk tidak menyerah serta tetap memilih untuk bangkit di tengah segala tekanan, tantangan, dan keterpurukan hingga akhirnya mampu menyelesaikan seluruh rangkaian perjalanan akademis ini.
    \item Orang tua dan keluarga tercinta yang senantiasa memberikan doa yang tiada putus, semangat, kasih sayang, serta dukungan moral maupun material sepanjang masa perkuliahan hingga selesainya Tugas Akhir ini.
    \item Bapak Dr. rer. nat. Bintoro Anang Subagyo, selaku dosen pembimbing yang telah dengan penuh kesabaran meluangkan waktu, tenaga, pikiran, serta memberikan arahan dan ilmu yang sangat berharga dalam membimbing penulis menyelesaikan Tugas Akhir ini.
    \item Bapak dan Ibu dosen penguji atas kritik, saran, dan masukan konstruktif yang sangat membantu penulis dalam menyempurnakan Tugas Akhir ini.
    \item Bapak Prof. Agus Purwanto, selaku dosen yang senantiasa memberikan motivasi, inspirasi, dan dorongan semangat bagi penulis untuk menjadi seorang fisikawan yang berintegritas.
    \item Bapak Dr. Heru Sukamto, M.Si., selaku dosen yang telah memberikan banyak ilmu, teladan, dan pembelajaran berharga kepada penulis selama masa studi.
    \item Bapak Dr. Muhammad Taufiqi, selaku dosen yang telah memberikan dukungan luar biasa serta membuka jalan bagi penulis untuk mendapatkan pengalaman akademik berharga ke luar negeri untuk pertama kalinya.
    \item Bapak Dr. Lila Yuwana, selaku dosen wali yang senantiasa memberikan bimbingan, masukan bijak, serta membantu penulis dalam mengarahkan alur studi selama berkuliah.
    \item Seluruh dosen dan staf administrasi maupun sarana-prasarana Departemen Fisika FSAD ITS yang telah membagikan ilmu yang bermanfaat serta menyediakan fasilitas pendukung selama proses perkuliahan.
    \item Teman-teman LaFTiFA 2022, khususnya Alief, Alva, Ilham, dan Naufal, yang selalu setia menemani, memberikan hiburan, serta menjadi tempat bertukar cerita dan keluh kesah di kala penulis mengerjakan Tugas Akhir ini.
    \item Teman-teman RUDAL Corp., atas rasa kekeluargaan, solidaritas yang erat, serta dukungan yang tak henti-hentinya kepada penulis sepanjang berkuliah di Fisika ITS.
    \item Teman-teman seperjuangan angkatan Fisika ITS 2022 yang telah mewarnai hari-hari perkuliahan dengan ruang diskusi, motivasi, dan kebersamaan yang tidak akan terlupakan.
    \item Seluruh pihak yang tidak dapat disebutkan satu per satu, yang telah memberikan bantuan, doa, maupun dukungan dalam bentuk apa pun demi kelancaran penyelesaian Tugas Akhir ini.
\end{enumerate}
\indent Penulis menyadari bahwa tugas akhir ini masih jauh dari kesempurnaan. Oleh karena itu, penulis mengharapkan kritik dan saran yang bersifat membangun. Semoga tugas akhir ini dapat memberikan manfaat bagi pengembangan ilmu pengetahuan.

\vspace{1cm}
\begin{flushright}
	Surabaya, Agustus 2026
	
	Penulis
\end{flushright}

\newpage
\thispagestyle{empty}
\vspace*{\fill}
	\begin{center}
		Halaman ini sengaja dikosongkan
	\end{center}
\vspace*{\fill}
\pagebreak
